\chapter{Network-behavior co-evolution}\label{ch:coevolution}
\begingroup\raggedright\sloppy

A network-only model treats the actors' attributes as fixed. In many applications the attributes change too, and they change \emph{because of} the network: adolescents drink more when their friends drink (influence), and they choose friends who drink like they do (selection). \nwsaom{} models both processes at once by adding a second dependent variable, an ordinal \emph{behavior}, that evolves between the same observed waves as the network. This is RSiena's network-behavior co-evolution model \cite{steglich2010}. Throughout the chapter the data are the 50 pupils of the Glasgow friendship study that ship with the package (\texttt{nwwebuse glasgow}): three friendship waves \texttt{glasgow1}--\texttt{glasgow3}, an alcohol variable \texttt{alcohol1}--\texttt{alcohol3} (five categories), smoking \texttt{smoke1}--\texttt{smoke3} (three categories) and sport \texttt{sport1}--\texttt{sport3}. These are the same data as RSiena's \texttt{s501}--\texttt{s503}, \texttt{s50a}, \texttt{s50s}; the node order differs, the results do not.

\section{The model in brief}

A behavior is a bounded integer variable, given as one Stata variable per wave. Between two waves the simulated process has two kinds of ministeps:
\begin{itemize}[nosep]
\item a \emph{network ministep}: an actor reconsiders one outgoing tie, exactly as in Chapter~\ref{ch:model};
\item a \emph{behavior ministep}: an actor changes its own behavior value by $-1$, $0$ or $+1$ (clamped to the observed minimum and maximum). The three alternatives are chosen with the same multinomial-logit rule as a network ministep, now over three alternatives instead of $n$.
\end{itemize}
One pooled waiting time is drawn from the sum of the network rate and the behavior rate. Which variable acts is chosen in proportion to its share of that total, and then the actor is chosen uniformly. This is RSiena's construction (verified against its \texttt{EpochSimulation.cpp} according to the help file). Every period has its own network rate and its own behavior rate.

Co-evolution models are always estimated by \emph{unconditional} Method of Moments, as in RSiena: both kinds of rates are estimated jointly with the effects, and they come with standard errors. The statistic of a network rate is the number of dyads in which the simulated end network differs from the start network; the statistic of a behavior rate is the sum of the absolute changes of the behavior. Statistics that link the two variables are \emph{lagged} as in RSiena: network effects that read the behavior (\opt{behsim}, \opt{behegox}, \dots) use the behavior at the \emph{start} of the period, behavior effects that read the network (\opt{avalt}, \dots) use the network at the start of the period. The ministep change statistics always use the current simulated state.

\section{A first model}

The behavior is declared with \opt{behavior()}, which takes the wave variables in temporal order. At least one baseline behavior effect is required, as \opt{outdegree} is on the network side: \opt{linear} (RSiena's \texttt{linear} shape).

\begin{lstlisting}
nwwebuse glasgow, nwclear
nwsaom, waves(glasgow1 glasgow2 glasgow3) outdegree reciprocity transtrip behsim ///
    behavior(alcohol1 alcohol2 alcohol3) linear quadratic avalt seed(12345) k3(300)
\end{lstlisting}
\begin{lstlisting}
SAOM co-evolution (Method of Moments), waves: glasgow1 glasgow2 glasgow3
Actors: 50                             Periods: 2
Behavior: alcohol1 alcohol2 alcohol3
-------------------------------------------------------------------------------
glasgow1_to~3 | Coefficient  Std. err.      z    P>|z|
--------------+----------------------------------------------------------------
    outdegree |  -2.760383   .1418124   -19.47   0.000
  reciprocity |   2.367478    .207555    11.41   0.000
    transtrip |   .6150662   .0746171     8.24   0.000
       behsim |   1.446294   .8003719     1.81   0.071
   beh_linear |   .3965896   .2153328     1.84   0.066
beh_quadratic |   -.607822   .3080686    -1.97   0.048
    beh_avalt |   1.317363   .7485551     1.76   0.078
-------------------------------------------------------------------------------
             |   period1    period2
network_rate |    6.4518     5.1274
behavior_rate|    1.3194     1.8041
Overall maximum convergence ratio:  0.222
\end{lstlisting}
(Rate standard errors and the convergence $t$-ratios are omitted here; \texttt{e(rates)}, \texttt{e(rates\_beh)} and \texttt{e(tconv)} hold them.) The first four coefficients belong to the network side, the last three to the behavior side. Network coefficients keep their plain names; behavior coefficients start with \texttt{beh\_}. The network effects use the \opt{waves()} and \opt{seed()} options exactly as in Chapter~\ref{ch:start}.

The help file documents a comparison with RSiena 1.6.6 on this model (RSiena estimate, standard error in brackets; the \nwsaom{} column is the value in the help file, the run above with \opt{k3(300)} gives values of the same size, e.g.\ \opt{behsim} 1.446, \texttt{beh\_avalt} 1.317):

\begin{center}
\begin{tabular}{lrr}
\toprule
effect & RSiena & \nwsaom{} (help file) \\
\midrule
friendship rate, period 1 / 2 & 6.48 / 5.17 & 6.43 / 5.12 \\
outdegree & $-2.760$ (.145) & $-2.752$ (.188) \\
reciprocity & 2.354 (.198) & 2.364 (.217) \\
transitive triplets & 0.617 (.077) & 0.615 (.095) \\
drinking similarity (\opt{behsim}) & 1.484 (.626) & 1.431 (.842) \\
drinking rate, period 1 / 2 & 1.31 / 1.81 & 1.31 / 1.73 \\
linear shape & 0.392 (.205) & 0.386 (.216) \\
quadratic shape & $-0.589$ (.311) & $-0.567$ (.391) \\
average alter & 1.286 (.774) & 1.223 (.882) \\
\bottomrule
\end{tabular}
\end{center}

The quantities returned by the fit that are specific to co-evolution are \texttt{e(rates\_beh)} and \texttt{e(rates\_beh\_se)} (the behavior rates; with several behaviors one row per behavior), \texttt{e(has\_behavior)}, and \texttt{e(engine)}, which says whether the native plugin or the Mata engine did the simulations.

\section{Shape effects}

Shape effects describe the behavior's own distribution, independent of the network. They are the behavior-side analogue of \opt{outdegree}.

\begin{longtable}{@{}p{3.4cm}p{3.2cm}p{7.6cm}@{}}
\toprule
option & RSiena effect & what it does \\
\midrule
\endhead
\opt{linear} & \texttt{linear} & Baseline tendency to higher values. Ministep statistic: the raw change ($\pm1$); observed statistic: sum of the current values. Required unless \opt{behshape()} replaces the shape effects. \\
\opt{quadratic} & \texttt{quad} & Quadratic shape. Copied from RSiena including its quirk: the ministep uses the \emph{centered} value ($(2(z-\bar z)+d)\,d$), while the global statistic sums the raw, uncentered squared value. \\
\opt{behshape(quad\_cc)}, \opt{behshape(quad\_nc)} & \texttt{quad} variants (centred / non-centred) & Alternatives to \opt{quadratic}; use them \emph{instead of} it, never together. \\
\opt{behshape(threshold:\#)} & \texttt{threshold} & Indicator for values at or above a threshold. \\
\opt{behshape(simallnear:\#)}, \opt{behshape(simallfar:\#)} & \texttt{simAllNear}, \texttt{simAllFar} & Similarity to \emph{all} actors within / beyond a distance (defaults 2 and 4). \\
\opt{behshape(avgroup[:\#])}, \opt{behshape(totgroup[:\#])} & \texttt{avGroup}, \texttt{totGroup} & Group-average and group-total effects (default parameter 1). \\
\opt{behshape(constant)} & \texttt{constant} & Coded but not estimable: its statistic is minus the behavior rate statistic, so the Jacobian is singular. \\
\bottomrule
\end{longtable}

Coefficients of \opt{behshape()} effects are named \texttt{beh\_}\emph{effect}[\texttt{(}\emph{\#}\texttt{)}]. The help file reports agreement with \texttt{siena07} on the s50 data (for example \opt{quad\_cc}, linear 1.791 / quad $-0.448$ / avAlt 1.057 against 1.790 / $-0.450$ / 1.071).

\section{Influence effects}

Influence effects make an actor's behavior depend on its network neighbours.

\subsection{The two main effects: avalt and avsim}

\opt{avalt} (RSiena's \texttt{avAlt}) pulls an actor's value towards the average value of its out-neighbours: the ministep contribution is $d\,(\bar z_{\mathrm{alters}} - \bar z)$ (zero for an actor without out-ties), with $\bar z$ the overall behavior mean. \opt{avsim} (RSiena's \texttt{avSim}) is the alternative parameterisation: the actor maximises its average \emph{similarity} $1 - |z_i - z_j|/\text{range}$ to its neighbours, net of a data-derived centring constant (the similarity mean, computed from the observed data and never supplied by the user). One RSiena quirk is copied: this constant is exactly 0 when the pooled behavior has zero variance. The two effects are alternatives; the examples below use them separately.

\begin{lstlisting}
nwsaom, waves(glasgow1 glasgow2 glasgow3) outdegree reciprocity transtrip ///
    behegox behaltx behavior(alcohol1 alcohol2 alcohol3) linear quadratic avsim ///
    seed(1) k3(300)
\end{lstlisting}
\begin{lstlisting}
    outdegree |  -2.647354    .140056   -18.90   0.000
  reciprocity |   2.429259   .2282148    10.64   0.000
    transtrip |   .6047165   .0743783     8.13   0.000
      behegox |   .0447982   .0951204     0.47   0.638
      behaltx |  -.0463469   .0847807    -0.55   0.585
   beh_linear |   .3868733   .1699837     2.28   0.023
beh_quadratic |  -.0535451   .1064301    -0.50   0.615
    beh_avsim |   4.123988    2.13372     1.93   0.053
\end{lstlisting}

\subsection{Endowment and creation roles}

Every shape and influence effect has the three RSiena roles: evaluation (the default, direction-blind), endowment (counts only when the value decreases) and creation (counts only when it increases). The split is requested with the pair \opt{quadraticendow}/\opt{quadraticcreation}, \opt{avaltendow}/\opt{avaltcreation} or \opt{avsimendow}/\opt{avsimcreation} (equivalently \opt{endow(avalt) creation(avalt)}); both members must be given and the plain effect must be absent, because all three together are exactly collinear. Splits are weakly identified in general:

\begin{lstlisting}
nwsaom, waves(glasgow1 glasgow2 glasgow3) outdegree reciprocity transtrip behsim ///
    behavior(alcohol1 alcohol2 alcohol3) linear quadratic avaltendow avaltcreation ///
    seed(1) k3(300)
\end{lstlisting}
\begin{lstlisting}
 beh_avalt_endow |  -2.141754   3.210858    -0.67   0.505
beh_avalt_creation|   3.065424   4.502735     0.68   0.496
\end{lstlisting}
Here the fit converges (maximum convergence ratio 0.303, above the 0.25 guideline) but the two roles have standard errors of 3 to 4.5; the data do not separate them. Such fits can also stop with the \texttt{thetaBound} error or a near-singular covariance matrix; both mean ``use the plain effect''. The split of \opt{linear} (\opt{linearendow}/\opt{linearcreation}) is refused: the sums of increases and decreases add to the behavior-rate statistic, so with an estimated behavior rate the model is not identified (RSiena also reports a singular covariance matrix).

\subsection{Network effects on behavior: beheff, behnet2 and friends}

RSiena offers many further behavior effects that depend on the network. All are native-only.

\begin{longtable}{@{}p{3.6cm}p{9.2cm}@{}}
\toprule
option & what it adds \\
\midrule
\endhead
\opt{beheff(tokens)} & Ego behavior $\times$ position or alter summaries. \emph{Position:} \texttt{indeg outdeg recipdeg isolate outisolate}. \emph{Alter summaries:} \texttt{totalt avaltpop totaltpop maxalt minalt avrecalt totrecalt avinalt totinalt popalt totpopalt}. \emph{Similarity families:} \texttt{totsim avatthigher avattlower totatthigher totattlower avsimpopalt totsimpopalt avsimpopego avinsim totinsim avinsimpopalt totinsimpopalt avsimrecip totsimrecip avsimrecpop totsimrecpop}. \emph{Distance two:} \texttt{avaltdist2 totaltdist2 avtaltdist2 totaaltdist2 avinaltdist2 totinaltdist2 avtinaltdist2 totainaltdist2}. The tokens are the lower-cased RSiena short names (\texttt{totalt} = \texttt{totAlt}, \texttt{avaltpop} = \texttt{avAltPop}, \dots). Coefficients \texttt{beh\_}\emph{token}. \\
\opt{behnet2(effect[:\#])} & Further effects: \texttt{indegavgroup indegtotgroup avaltcc totaltcc avinaltdist2nc totinaltdist2nc avtinaltdist2nc totainaltdist2nc avinsimdist2 totinsimdist2 totgwdspffalt totgwdspfbalt totgwdspffaltnc totgwdspfbaltnc varalt avsimvaralt behdensetriads simdensetriads}. The number after a colon is RSiena's effect parameter (group effects: centring value, default 1; gwdsp: weight, default 69; dense triads: 5 or 6). \\
\opt{behfrom(varlist)} & Covariate main effect on the behavior (RSiena \texttt{effFrom}): behavior tends to be higher where the covariate is higher; contribution is the change times the centred covariate. Constant or \opt{tvcov()} covariates (stem name), first behavior, at most 8; coefficient \texttt{beh\_from\_}\emph{var}. \\
\opt{behavxalt()}, \opt{behavsamexalt()}, \opt{behtotxalt()}, \opt{behtotsamexalt()} & Covariate-moderated influence (RSiena \texttt{avXAlt}, \texttt{avSameXAlt}, \texttt{totXAlt}, \texttt{totSameXAlt}). \\
\opt{behx(family:var)} & Behavior effects of the network $\times$ a covariate, 32 families (\texttt{avxinalt totxinalt avxrecalt totxrecalt avsamexinalt totsamexinalt avsamexrecalt totsamexrecalt avaltsamex totaltsamex avrecaltsamex totrecaltsamex avaltaltx avsimaltx totsimaltx minxalt maxxalt degabscontrx degposcontrx degnegcontrx} and the distance-two and ego families \texttt{avxaltdist2 \dots{} totaltegox}). \emph{var} is a numeric variable, a \opt{tvcov()} stem, or \texttt{beh2}, \texttt{beh3}, \dots{} (another co-evolving behavior, Section~\ref{sec:multibeh}). Coefficients \texttt{beh\_}\emph{family}\texttt{\_}\emph{var}. \\
\opt{behw(family:matrix)} & Behavior effects of the network $\times$ a dyadic covariate $W$: \texttt{avaltw totaltw avinaltw totinaltw avwalt totwalt avwinalt totwinalt avsimw totsimw}. \emph{matrix} is a Stata $n\times n$ matrix, centred like \opt{dyadcov()}; at most 8 matrices. \opt{behwparam(1|2)} chooses RSiena's parameter of \texttt{avAltW}, \texttt{avInAltW}, \texttt{avSimW} (2, the default, divides by the sum of the weights; 1 by the number of alters; names then carry \texttt{\_p1}). \\
\opt{tvbehw(family:stem)} & The same ten effects with a changing dyadic covariate; matrices \emph{stem}1 \dots{} \emph{stem}($W-1$), one per period. May be mixed with \opt{behw()}. \\
\bottomrule
\end{longtable}

The help file states that for all of these the observed statistics equal \texttt{siena07}'s targets and the estimates agree within simulation noise, with a handful of exceptions that are unstable or degenerate in \texttt{siena07} as well (for example \texttt{avAltW} and \texttt{avSimW} with parameter 2 when the weights nearly sum to zero). Combining many weakly identified effects makes the standard errors explode; the next fit is shown only to illustrate the syntax:

\begin{lstlisting}
nwsaom, waves(glasgow1 glasgow2 glasgow3) outdegree reciprocity transtrip behsim ///
    behavior(alcohol1 alcohol2 alcohol3) linear quadratic avalt ///
    beheff(totalt indeg) behfrom(sport1) seed(1) k3(300)
\end{lstlisting}
\begin{lstlisting}
   beh_linear |    .126359   .5581673     0.23   0.821
beh_quadratic |  -.5475244   .3459674    -1.58   0.114
    beh_avalt |   .1023272   2.435999     0.04   0.966
   beh_totalt |    .382001   .7636128     0.50   0.617
    beh_indeg |   .1098258   .2436868     0.45   0.652
beh_from_sport1|   .2236675   .3975685     0.56   0.574
\end{lstlisting}
(\texttt{avalt} and \texttt{totalt} are nearly collinear in these data; the standard error of \texttt{beh\_avalt} grows from 0.75 in the first model to 2.4.)

An example with a dyadic covariate: the first wave of the friendship network used as the weight matrix $W$.
\begin{lstlisting}
nwtomata glasgow1, mat(M1)
mata: st_matrix("DM", M1)
nwsaom, waves(glasgow1 glasgow2 glasgow3) outdegree reciprocity ///
    behavior(alcohol1 alcohol2 alcohol3 : linear) behw(avsimw:DM) seed(1) k3(300)
\end{lstlisting}
\begin{lstlisting}
    outdegree |  -2.356908   .1159083   -20.33   0.000
  reciprocity |   2.822803   .1968733    14.34   0.000
   beh_linear |   .3552337   .1554802     2.28   0.022
beh_avsimw_DM |   3.467163   1.336423     2.59   0.009
\end{lstlisting}
(This is only a demonstration of the syntax: using the network itself as a dyadic covariate is not a sensible substantive model.)

\section{Selection effects: the behavior shapes the network}\label{sec:selection}

Selection effects are \emph{network} effects whose covariate is the dependent behavior. Unlike \opt{simcov()} and \opt{nodeocov()}, which read a fixed covariate, they read the behavior as it changes during the simulation. In the ministeps they use the live value; in the observed and simulated statistics the start-of-period value (lagging, see above). Their coefficients are listed among the network coefficients, after the other network effects.

\begin{longtable}{@{}p{3.4cm}p{3.0cm}p{7.8cm}@{}}
\toprule
option & RSiena effect & what it does \\
\midrule
\endhead
\opt{behsim} & \texttt{simX} (dependent behavior) & Actors prefer ties to others with similar behavior: tie contribution $1 - |z_i - z_j|/\text{range} - \text{simMean}$, with range and similarity mean the behavior's own (the same constants \opt{avsim} uses). The centring matches RSiena, so the outdegree coefficient is comparable to RSiena's. \\
\opt{behegox} & \texttt{egoX} & Sender's current behavior minus the overall behavior mean ("alcohol ego"). \\
\opt{behaltx} & \texttt{altX} & Receiver's behavior minus the mean ("alcohol alter"). \\
\opt{behegoplusalt} & \texttt{egoPlusAltX} & Sum of both. \\
\opt{netbeh(effects)} & the covariate effects with the dependent behavior as covariate & Any of \texttt{samex diffx higher egosqx altsqx samexrecip absdiffx egoxaltx samextranstrip diffxtranstrip homxtranstrip samexoutact diffxoutact samexreciact diffxreciact homxoutact altxoutact inpopx outpopx inactx outactx samexinpop diffxinpop egoplusaltsqx diffsqx egodiffx althigheregox unequalx egolthresholdx egorthresholdx altlthresholdx altrthresholdx degabsdiffx degposdiffx degnegdiffx}, plus the similarity versions \texttt{simdist2 simegodist2 simegoindist2}. The behavior is centred by its overall mean in the simulation and taken at the start of the period in the statistic. Coefficients \texttt{bv\_}\emph{effect}. \\
\opt{netx(effect:var)} with \opt{netbeh()} & distance-two and covariate-closure effects & The \opt{netx()} effects of Chapter~\ref{ch:effects} also work with the behavior as covariate when given via \opt{netbeh()} (all except \texttt{homxoutact2}). \\
\bottomrule
\end{longtable}

\begin{lstlisting}
nwsaom, waves(glasgow1 glasgow2 glasgow3) outdegree reciprocity transtrip ///
    netbeh(samex egosqx) behavior(alcohol1 alcohol2 alcohol3) linear quadratic avalt ///
    seed(1) k3(300)
\end{lstlisting}
\begin{lstlisting}
    outdegree |    -3.0174   .2388013   -12.64   0.000
  reciprocity |   2.391431   .2237419    10.69   0.000
    transtrip |   .6101854   .0867924     7.03   0.000
     bv_samex |   .6676742   .3624426     1.84   0.065
    bv_egosqx |   .1200029    .090639     1.32   0.186
   beh_linear |   .3999268   .2106728     1.90   0.058
beh_quadratic |  -.5643059   .3146445    -1.79   0.073
    beh_avalt |     1.2566   .7310308     1.72   0.086
\end{lstlisting}

According to the help file and decision log, 32 \opt{netbeh()} estimates agree with \texttt{siena07} on the s50 data; \texttt{degNegDiffX} is degenerate in both programs.

\section{Several behaviors}\label{sec:multibeh}

Any number of behaviors can co-evolve with the network. In \opt{behavior()} the behaviors are separated by \texttt{|}, and each can carry its own effect list after a colon; without a colon the global effect options apply to every behavior (one behavior with a colon is identical to the global form).

\begin{lstlisting}
nwsaom, waves(glasgow1 glasgow2 glasgow3) outdegree reciprocity transtrip ///
    behavior(alcohol1 alcohol2 alcohol3 : linear quadratic avalt behsim | ///
             smoke1 smoke2 smoke3 : linear quadratic avalt) seed(1) k3(300)
\end{lstlisting}
\begin{lstlisting}
SAOM co-evolution of the network and 2 behaviors (Method of Moments)
            outdegree |  -2.754319   .1429746   -19.26   0.000
          reciprocity |   2.357028    .212744    11.08   0.000
            transtrip |   .6164275   .0827658     7.45   0.000
               behsim |   1.430266   .7080562     2.02   0.043
   beh_alcohol_linear |   .3934624   .2099502     1.87   0.061
beh_alcohol_quadratic |  -.5937406   .3414086    -1.74   0.082
    beh_alcohol_avalt |   1.294857   .8745685     1.48   0.139
     beh_smoke_linear |  -1.862156   .3160527    -5.89   0.000
  beh_smoke_quadratic |   1.947808   .4002872     4.87   0.000
      beh_smoke_avalt |   1.761282    .858102     2.05   0.040
\end{lstlisting}
\opt{behsim} is written inside the group of the \emph{first} behavior. (In the version tested for this manual, a global \opt{behsim} given outside the groups together with several behaviors was accepted without error but produced no \texttt{behsim} coefficient and the same estimates as the model without it, so always put it inside the first group.)

Coefficients are named \texttt{beh\_}\emph{name}\texttt{\_}\emph{effect}, with \emph{name} the first variable without trailing digits. \texttt{e(rates\_beh)} has one row per behavior (values of the model above):
\begin{lstlisting}
                period1    period2
rate_alcohol  1.3350     1.7769
  rate_smoke  2.9479     2.6810
\end{lstlisting}
The help file's validation against RSiena 1.6.6 (alcohol and smoking, \texttt{s501}--\texttt{s503}) gives, for example, alcohol linear 0.42 (RSiena 0.43), smoking linear $-1.91$ to $-1.94$ ($-1.92$ to $-1.94$), smoking avAlt 1.92--1.97 (2.01--2.06), with the network and behavior rates agreeing.

The behaviors interact through the network, and in one further way: another behavior can be the \emph{covariate} of the behavior effects that take one, by writing \texttt{beh2}, \texttt{beh3}, \dots{} in place of a variable name in \opt{behfrom()}, \opt{behx()}, \opt{behavxalt()} and relatives (the statistics read its start-of-period values, the ministeps its current ones):
\begin{lstlisting}
nwsaom, waves(glasgow1 glasgow2 glasgow3) outdegree reciprocity transtrip ///
    behavior(alcohol1 alcohol2 alcohol3 : linear quadratic avalt | ///
             smoke1 smoke2 smoke3 : linear quadratic avalt) behfrom(beh2) seed(1) k3(300)
\end{lstlisting}
\begin{lstlisting}
beh_alcohol_from_beh2 |  -.7376024   1.023347    -0.72   0.471
\end{lstlisting}
(smoking as a main effect on drinking; not significant in these data.)

\section{Co-evolution with other features}

\begin{itemize}[nosep]
\item \textbf{Non-directed networks.} A symmetric network can co-evolve with a behavior; see Chapter~\ref{ch:networktypes}. It requires the native plugin with the threaded batch path.
\item \textbf{Two-mode networks.} Actors of a two-mode network may carry a behavior; see Chapter~\ref{ch:networktypes}.
\item \textbf{Composition change.} \opt{joinat()}/\opt{leaveat()} work with one or several behaviors (not with \opt{behsim} or endowment/creation effects), Chapter~\ref{ch:rates}.
\item \textbf{Time-varying covariates.} \opt{tvcov()} works with \opt{behavior()} for the network covariate effects; \opt{behfrom()}, \opt{behx()} and \opt{behavxalt()} accept a \opt{tvcov()} stem.
\item \textbf{Starting values.} \opt{behtheta0()} gives starting values for the behavior-side parameters.
\end{itemize}

\section{Estimation notes}

\begin{itemize}[nosep]
\item With few actors, \opt{avalt} can make the joint Robbins--Monro fit diverge: a stronger pull makes the ministep more deterministic, which makes the pull look stronger. The help file reports that this is a small-sample identification problem (checked: the Jacobian is well conditioned and the simulator unbiased at the true parameter), resolved by using more actors or more behavior activity. A divergence stops with an explicit \texttt{thetaBound} error (return code 498), as in RSiena.
\item Several of the examples above have a maximum convergence ratio between 0.22 and 0.33; RSiena's guideline is 0.25. Rerun with another \opt{seed()} or a larger \opt{k3()} when a fit is above the limit and the reported $t$-ratios of individual parameters exceed 0.1.
\item Random or weakly varying behaviors are usually not identified; the test suite therefore accepts return codes 0 and 498 for synthetic co-evolution data.
\end{itemize}

\section{Limits}\label{sec:coev-limits}

The following are things RSiena can do (or that this chapter would naturally want) and \nwsaom{} cannot, taken from the help file, the remarks, the decision log, and checks run for this chapter (``checked'' means an error was observed).

\begin{itemize}
\item \textbf{Behavior type.} Behaviors are bounded integers, one Stata variable per wave. There is no continuous behavior.
\item \textbf{Linear endowment/creation is refused} (not identified with an estimated behavior rate; RSiena also yields a singular covariance matrix). Network \opt{outdegreeendow} with \opt{outdegreecreation} is refused too, with or without a behavior. The other endow/creation splits of behavior effects are allowed but weakly identified.
\item \textbf{Several behaviors.} The following are not combinable with several behaviors: \opt{symmetric} (checked), \opt{present()} (checked; use \opt{joinat()}/\opt{leaveat()} with integer times), \opt{ratecov()}, \opt{missnet()}, \opt{structural()}, network \opt{endow()}/\opt{creation()}, and \texttt{estat gof}/\texttt{estat mems}. \opt{behsim} refers to the first behavior only (checked: \texttt{behsim} inside the group of the second behavior is rejected with return code 198; a global \opt{behsim} with several behaviors is silently ignored, see Section~\ref{sec:multibeh}). The same first-behavior restriction holds for \opt{behegox}, \opt{behaltx}, \opt{behegoplusalt} and \opt{netbeh()}.
\item \textbf{Cross-behavior effects} exist only through the covariate slot (\texttt{beh2} in \opt{behfrom()}, \opt{behx()}, \opt{behavxalt()} \dots). No other effects between two dependent behaviors are offered.
\item \textbf{\opt{ratecov()} with \opt{behavior()}} is refused (checked: ``ratecov() is not yet supported for co-evolution models''). Structural rate effects are also refused with a behavior (Chapter~\ref{ch:rates}).
\item \textbf{\opt{missbeh()} with several behaviors} is not available in the form the design note describes: a call \texttt{missbeh(a1 a2 a3 | b1 b2 b3)} is rejected (checked). Note also that \opt{missbeh()} variables are 0/1 flags, not Stata missing values.
\item \textbf{Native plugin only.} Since the project policy of 2026-10-22 new effects have no Mata fallback. Almost everything on this page after \opt{avsim} (\opt{beheff()} tokens beyond the first 20, \opt{behnet2()}, \opt{behx()}, \opt{behw()}, \opt{tvbehw()}, \opt{netbeh()}, \opt{behshape()}, \opt{behavxalt()} and relatives) needs the C plugin. The decision log notes that the Linux and Windows plugins lagged behind the macOS one in protocol versions; whether current builds exist for those platforms was not verified.
\item \textbf{Degenerate or unestimable effects.} \opt{behshape(constant)} cannot be estimated. For a number of effects \texttt{siena07} itself is unstable or degenerate (\opt{beheff()}: \texttt{outisolate totaltpop maxalt avaltpop avattlower totattlower totsim}; \opt{behw()} with parameter 2; some \opt{netbeh()} effects), so no reference estimate exists. In a non-directed network \texttt{degabsdiffx}, \texttt{degposdiffx} and \texttt{higher} with the behavior as covariate end in a singular phase-3 Jacobian (return code 505) whereas \texttt{siena07} estimates them; \texttt{minxalt}/\texttt{maxxalt} with \opt{symmetric} stay open.
\item \textbf{Symmetric co-evolution} is only partly validated against RSiena, because RSiena's own fit of the model is unstable on the test data. The agreement is by simulation from known parameters.
\item \textbf{Behavior inside \texttt{nwsaom multinet}} (several networks) is not available; see Chapter~\ref{ch:networktypes}.
\item \textbf{Inference.} Estimation is Method of Moments only. Maximum-likelihood and Bayesian SAOM estimation (\texttt{siena07} with \texttt{maxlike}, \texttt{sienaBayes}) are not planned.
\end{itemize}
\endgroup
